\section{Versiones para $CW$-complejos}

% ------------------------------------------------------------------------
\subsection{Subcomplejos}

\begin{frame}[allowframebreaks]
  \frametitle{Subespacios con buenos entornos}

  En esta sección, $A,B\subseteq X$ son subespacios con $X=A\cup B$ para los
  que existen abiertos $U\supseteq A$ y $V\supseteq B$ tales que
  \[
    A\subseteq U,\qquad B\subseteq V,\qquad A\cap B\subseteq U\cap V
  \]
  son retractos por deformación. Nótese que $X\supseteq U\cup V\supseteq A\cup B=X$,
  luego $\{U,V\}$ es un recubrimiento de $X$.

  \framebreak

  El caso principal es el de los $CW$-complejos.

  \begin{proposition}[Hatcher, Prop.~A.5]
    Sea $X$ un $CW$-complejo. A cada subcomplejo $A\subseteq X$ se le puede
    asociar un entorno abierto $N(A)\subseteq X$ de modo que
    \begin{enumerate}
      \item $A$ es un retracto por deformación de $N(A)$;
      \item $N(A)\cap N(B)=N(A\cap B)$ para cualesquiera subcomplejos $A,B$.
    \end{enumerate}
  \end{proposition}

  Por tanto, si $X$ es un $CW$-complejo y
  $A,B\subseteq X$ son subcomplejos con $X=A\cup B$, los abiertos $U=N(A)$
  y $V=N(B)$ cumplen la hipótesis, ya que $U\cap V=N(A\cap B)$.
\end{frame}

\begin{frame}[allowframebreaks]
  \frametitle{Reemplazar $A$ y $B$ por $U$ y $V$}

  Como los retractos de deformación son equivalencias homotópicas, el siguiente resultado es inmediato.

  \begin{lemma}
    Sean $A,B\subseteq X$ con $X=A\cup B$, y $U\supseteq A$, $V\supseteq B$
    abiertos como antes. Las inclusiones inducen isomorfismos
    \begin{gather*}
      H_n(A)\cong H_n(U),\qquad H_n(B)\cong H_n(V),\qquad
      H_n(A\cap B)\cong H_n(U\cap V),\\
      H_n(X,A)\cong H_n(X,U),\qquad H_n(B,A\cap B)\cong H_n(V,U\cap V),
    \end{gather*}
    y lo mismo en homología reducida.
  \end{lemma}

\end{frame}

% ------------------------------------------------------------------------
\subsection{Escisión y Mayer--Vietoris}

\begin{frame}[allowframebreaks]
  \frametitle{Escisión}

  \begin{theorem}
    Sean $A,B\subseteq X$ como antes; por ejemplo, subcomplejos de un
    $CW$-complejo con $X=A\cup B$. Entonces la inclusión de pares
    $(B,A\cap B)\hookrightarrow(X,A)$ induce isomorfismos
    \[
      H_n(B,A\cap B)\cong H_n(X,A)
    \]
    para todo $n$.
  \end{theorem}

  \framebreak

  \begin{proof}
    Las inclusiones de pares dan un diagrama conmutativo
    \[
      \begin{tikzcd}[ampersand replacement=\&]
        H_n(B,A\cap B)\arrow[r]\arrow{d}[swap]{\cong}
          \&H_n(X,A)\arrow{d}{\cong}\\
        H_n(V,U\cap V)\arrow{r}{\cong}\&H_n(X,U)
      \end{tikzcd}
    \]
    Las flechas verticales son isomorfismos por el lema anterior. La
    inferior lo es por el teorema de escisión aplicado al recubrimiento
    abierto $\{U,V\}$. Por tanto la
    superior también es un isomorfismo.
  \end{proof}
\end{frame}

\begin{frame}[allowframebreaks]
  \frametitle{Mayer--Vietoris}

  \begin{theorem}
    Sean $A,B\subseteq X$ como antes; por ejemplo, subcomplejos de un
    $CW$-complejo con $X=A\cup B$. Con las inclusiones $i_A,i_B,j_A,j_B$ de
    antes, existe una sucesión exacta larga
    \begin{multline*}
      \cdots\to H_n(A\cap B)\xrightarrow{\left(\begin{smallmatrix}{i_A}_*\\-{i_B}_*\end{smallmatrix}\right)}H_n(A)\oplus H_n(B)
      \xrightarrow{\left(\begin{smallmatrix}{j_A}_* & {j_B}_*\end{smallmatrix}\right)}H_n(X)\\
      \xrightarrow{\ \delta\ }H_{n-1}(A\cap B)\to\cdots\to H_0(X)\to0.
    \end{multline*}
    Si además $A\cap B\neq\varnothing$, lo mismo vale en homología reducida.
  \end{theorem}

  \framebreak

  \begin{proof}
    Como $\{U,V\}$ es un recubrimiento abierto de $X$, tenemos la sucesión
    de Mayer--Vietoris de $U$ y $V$. Por el
    lema anterior, las inclusiones inducen isomorfismos
    \[
      H_n(A\cap B)\cong H_n(U\cap V),\qquad
      H_n(A)\cong H_n(U),\qquad
      H_n(B)\cong H_n(V),
    \]
    también en homología reducida, 
    luego el teorema se sigue.
  \end{proof}
\end{frame}
